% GENERATED BY src/python/lmi/tables/tab03_employment.py -- DO NOT EDIT
%==============================================================================%
% tab03_employment.tex
%
% Purpose : Tabelle 3 der Replikation -- Wahrscheinlichkeit bezahlter Arbeit,
%           durchschnittliche marginale Effekte in Prozentpunkten. Entspricht
%           Table 3 in Kanas & Kosyakova, S. 20 (Bib-Key: Kanas:2023).
% Source  : SOEPremote job 244566, src/stata/remote/tab03_employment.do
%           Rohlisting: results/transcripts/src_stata_remote_tab03_employment.do.eml
%           Transkript: results/comparisons/tab03_employment.md
% Needs   : booktabs, tabularx (beide bereits in bachelorarbeit.sty geladen)
% Usage   : % GENERATED BY src/python/lmi/tables/tab03_employment.py -- DO NOT EDIT
%==============================================================================%
% tab03_employment.tex
%
% Purpose : Tabelle 3 der Replikation -- Wahrscheinlichkeit bezahlter Arbeit,
%           durchschnittliche marginale Effekte in Prozentpunkten. Entspricht
%           Table 3 in Kanas & Kosyakova, S. 20 (Bib-Key: Kanas:2023).
% Source  : SOEPremote job 244566, src/stata/remote/tab03_employment.do
%           Rohlisting: results/transcripts/src_stata_remote_tab03_employment.do.eml
%           Transkript: results/comparisons/tab03_employment.md
% Needs   : booktabs, tabularx (beide bereits in bachelorarbeit.sty geladen)
% Usage   : % GENERATED BY src/python/lmi/tables/tab03_employment.py -- DO NOT EDIT
%==============================================================================%
% tab03_employment.tex
%
% Purpose : Tabelle 3 der Replikation -- Wahrscheinlichkeit bezahlter Arbeit,
%           durchschnittliche marginale Effekte in Prozentpunkten. Entspricht
%           Table 3 in Kanas & Kosyakova, S. 20 (Bib-Key: Kanas:2023).
% Source  : SOEPremote job 244566, src/stata/remote/tab03_employment.do
%           Rohlisting: results/transcripts/src_stata_remote_tab03_employment.do.eml
%           Transkript: results/comparisons/tab03_employment.md
% Needs   : booktabs, tabularx (beide bereits in bachelorarbeit.sty geladen)
% Usage   : % GENERATED BY src/python/lmi/tables/tab03_employment.py -- DO NOT EDIT
%==============================================================================%
% tab03_employment.tex
%
% Purpose : Tabelle 3 der Replikation -- Wahrscheinlichkeit bezahlter Arbeit,
%           durchschnittliche marginale Effekte in Prozentpunkten. Entspricht
%           Table 3 in Kanas & Kosyakova, S. 20 (Bib-Key: Kanas:2023).
% Source  : SOEPremote job 244566, src/stata/remote/tab03_employment.do
%           Rohlisting: results/transcripts/src_stata_remote_tab03_employment.do.eml
%           Transkript: results/comparisons/tab03_employment.md
% Needs   : booktabs, tabularx (beide bereits in bachelorarbeit.sty geladen)
% Usage   : \input{tables/tab03_employment}
% Author  : Emre Suemer, 2026-08-13
%==============================================================================%

\makeatletter
\@ifundefined{NC@rewrite@Z}{%
  \newcolumntype{Z}{>{\raggedright\arraybackslash\hangindent=1.9em\hangafter=1}X}%
}{}
\makeatother

\begin{table}[htbp]
  \centering
  \caption{Wahrscheinlichkeit bezahlter Arbeit, marginale Effekte in Prozentpunkten}
  \label{tab:erwerbstaetigkeit}
  \footnotesize
  \setlength{\tabcolsep}{4pt}%
  \begin{tabularx}{\textwidth}{@{}Zrr@{}}
    \toprule
     & Modell 1.1 & Modell 1.2 \\
     & p.p. (SE) & p.p. (SE) \\
    \midrule
    \textbf{Kreisangebot an Integrationskursen, std.} & 2{,}98$^{*}$ (1{,}19) & 0{,}89 (1{,}78) \\
    Aufenthaltsdauer, in Monaten (/12) & 10{,}75$^{**}$ (1{,}60) & 10{,}99$^{**}$ (1{,}62) \\
    \quad $\times$ Kreisangebot an Integrationskursen &  & 0{,}81 (0{,}75) \\
    Bildungsjahre vor der Migration & 0{,}56$^{**}$ (0{,}19) & 0{,}56$^{**}$ (0{,}19) \\
    Erwerbserfahrung vor der Migration & 4{,}62$^{+}$ (2{,}39) & 4{,}57$^{+}$ (2{,}39) \\
    Weiblich & $-$15{,}83$^{**}$ (2{,}15) & $-$15{,}83$^{**}$ (2{,}15) \\
    Kinder unter 16 Jahren im Haushalt & $-$10{,}91$^{**}$ (2{,}15) & $-$10{,}90$^{**}$ (2{,}15) \\
    Einreise mit der Familie & 5{,}29$^{**}$ (1{,}96) & 5{,}33$^{**}$ (1{,}96) \\
    Unterst\"utzung durch Familie/Verwandte in DE & 4{,}10 (2{,}85) & 4{,}19 (2{,}85) \\
    Alter bei der Erstbefragung & $-$0{,}35$^{**}$ (0{,}09) & $-$0{,}35$^{**}$ (0{,}09) \\
    Traumatische Erfahrungen & $-$0{,}34 (2{,}31) & $-$0{,}36 (2{,}31) \\
    Status des Asylantrags (Ref.: keine Entscheidung) & & \\
    \quad -- Anerkannt & 5{,}20$^{*}$ (2{,}60) & 5{,}19$^{*}$ (2{,}60) \\
    \quad -- Abgelehnt & 8{,}24$^{**}$ (2{,}77) & 8{,}31$^{**}$ (2{,}76) \\
    Kreisanteil Ausl\"ander:innen, std. & $-$1{,}16 (2{,}79) & $-$1{,}12 (2{,}79) \\
    Kreisbev\"olkerungsdichte, std. & $-$0{,}89 (3{,}60) & $-$0{,}90 (3{,}59) \\
    Kreisarbeitslosenquote, std. & $-$4{,}45$^{+}$ (2{,}43) & $-$4{,}54$^{+}$ (2{,}44) \\
    Kreissorgen \"uber Zuwanderung, std. & 0{,}36 (1{,}20) & 0{,}34 (1{,}20) \\
    Kreisanteil CDU/CSU-W\"ahler:innen, std. & $-$1{,}92 (2{,}19) & $-$1{,}98 (2{,}20) \\
    Konstante & $-$2{,}63 (6{,}30) & $-$2{,}77 (6{,}29) \\
    \midrule
    \multicolumn{3}{@{}l}{\textit{Fixe Effekte}} \\
    \quad Erhebungsjahr & Ja & Ja \\
    \quad Erhebungsstichprobe & Ja & Ja \\
    \quad Bundesland & Ja & Ja \\
    \quad Herkunftsland (aggregiert) & Ja & Ja \\
    \midrule
    N Beobachtungen & 5.467 & 5.467 \\
    N Personen & 2.730 & 2.730 \\
    N Imputationen & 20 & 20 \\
    R\textsuperscript{2} (adj.) & 0{,}2135 & 0{,}2138 \\
    \bottomrule
  \end{tabularx}

  \vspace{0.75ex}
  \parbox{\textwidth}{\footnotesize
    \textit{Anmerkungen:} Eigene Berechnung, IAB-BAMF-SOEP-Befragung von
    Gefl\"uchteten 2016--2019, SOEP-Core v36 (Remote Edition), gewichtet.
    Lineare Wahrscheinlichkeitsmodelle; robuste, auf Personenebene geclusterte
    Standardfehler in Klammern. Bei standardisierten Variablen (std.) entspricht
    der Koeffizient dem Effekt einer Erh\"ohung um eine Standardabweichung.
    $^{+}$~p\,$<$\,0{,}1, $^{*}$~p\,$<$\,0{,}05, $^{**}$~p\,$<$\,0{,}01
    (zweiseitig). Alle 37 Koeffizienten- und Standardfehlerzellen stimmen im
    Signifikanzniveau mit \textcite[20]{Kanas:2023} \"uberein; die gr\"o\ss{}te
    Abweichung betr\"agt 0{,}28~p.p.
    \\[0.4ex]
    \textbf{Abweichung:} Die Kategorien des Asylstatus sind in
    der publizierten Tabelle falsch zugeordnet. Das Modell verwendet
    \texttt{ib1.asyl\_req\_stat}, und \texttt{00013\_data\_clean.do:1541}
    definiert Kategorie~1 als \emph{keine Entscheidung}; die Referenz ist somit
    \glqq keine Entscheidung\grqq{} und nicht \glqq anerkannt\grqq{}. Die
    Zuordnung oben ist die korrigierte und wurde gegen die Wertelabels der
    \texttt{mi estimate}-Ausgabe gepr\"uft. Inhaltlich bedeutet dies: sowohl
    anerkannte als auch abgelehnte Gefl\"uchtete arbeiten h\"aufiger als jene,
    deren Verfahren noch l\"auft.}
\end{table}

% Author  : Emre Suemer, 2026-08-13
%==============================================================================%

\makeatletter
\@ifundefined{NC@rewrite@Z}{%
  \newcolumntype{Z}{>{\raggedright\arraybackslash\hangindent=1.9em\hangafter=1}X}%
}{}
\makeatother

\begin{table}[htbp]
  \centering
  \caption{Wahrscheinlichkeit bezahlter Arbeit, marginale Effekte in Prozentpunkten}
  \label{tab:erwerbstaetigkeit}
  \footnotesize
  \setlength{\tabcolsep}{4pt}%
  \begin{tabularx}{\textwidth}{@{}Zrr@{}}
    \toprule
     & Modell 1.1 & Modell 1.2 \\
     & p.p. (SE) & p.p. (SE) \\
    \midrule
    \textbf{Kreisangebot an Integrationskursen, std.} & 2{,}98$^{*}$ (1{,}19) & 0{,}89 (1{,}78) \\
    Aufenthaltsdauer, in Monaten (/12) & 10{,}75$^{**}$ (1{,}60) & 10{,}99$^{**}$ (1{,}62) \\
    \quad $\times$ Kreisangebot an Integrationskursen &  & 0{,}81 (0{,}75) \\
    Bildungsjahre vor der Migration & 0{,}56$^{**}$ (0{,}19) & 0{,}56$^{**}$ (0{,}19) \\
    Erwerbserfahrung vor der Migration & 4{,}62$^{+}$ (2{,}39) & 4{,}57$^{+}$ (2{,}39) \\
    Weiblich & $-$15{,}83$^{**}$ (2{,}15) & $-$15{,}83$^{**}$ (2{,}15) \\
    Kinder unter 16 Jahren im Haushalt & $-$10{,}91$^{**}$ (2{,}15) & $-$10{,}90$^{**}$ (2{,}15) \\
    Einreise mit der Familie & 5{,}29$^{**}$ (1{,}96) & 5{,}33$^{**}$ (1{,}96) \\
    Unterst\"utzung durch Familie/Verwandte in DE & 4{,}10 (2{,}85) & 4{,}19 (2{,}85) \\
    Alter bei der Erstbefragung & $-$0{,}35$^{**}$ (0{,}09) & $-$0{,}35$^{**}$ (0{,}09) \\
    Traumatische Erfahrungen & $-$0{,}34 (2{,}31) & $-$0{,}36 (2{,}31) \\
    Status des Asylantrags (Ref.: keine Entscheidung) & & \\
    \quad -- Anerkannt & 5{,}20$^{*}$ (2{,}60) & 5{,}19$^{*}$ (2{,}60) \\
    \quad -- Abgelehnt & 8{,}24$^{**}$ (2{,}77) & 8{,}31$^{**}$ (2{,}76) \\
    Kreisanteil Ausl\"ander:innen, std. & $-$1{,}16 (2{,}79) & $-$1{,}12 (2{,}79) \\
    Kreisbev\"olkerungsdichte, std. & $-$0{,}89 (3{,}60) & $-$0{,}90 (3{,}59) \\
    Kreisarbeitslosenquote, std. & $-$4{,}45$^{+}$ (2{,}43) & $-$4{,}54$^{+}$ (2{,}44) \\
    Kreissorgen \"uber Zuwanderung, std. & 0{,}36 (1{,}20) & 0{,}34 (1{,}20) \\
    Kreisanteil CDU/CSU-W\"ahler:innen, std. & $-$1{,}92 (2{,}19) & $-$1{,}98 (2{,}20) \\
    Konstante & $-$2{,}63 (6{,}30) & $-$2{,}77 (6{,}29) \\
    \midrule
    \multicolumn{3}{@{}l}{\textit{Fixe Effekte}} \\
    \quad Erhebungsjahr & Ja & Ja \\
    \quad Erhebungsstichprobe & Ja & Ja \\
    \quad Bundesland & Ja & Ja \\
    \quad Herkunftsland (aggregiert) & Ja & Ja \\
    \midrule
    N Beobachtungen & 5.467 & 5.467 \\
    N Personen & 2.730 & 2.730 \\
    N Imputationen & 20 & 20 \\
    R\textsuperscript{2} (adj.) & 0{,}2135 & 0{,}2138 \\
    \bottomrule
  \end{tabularx}

  \vspace{0.75ex}
  \parbox{\textwidth}{\footnotesize
    \textit{Anmerkungen:} Eigene Berechnung, IAB-BAMF-SOEP-Befragung von
    Gefl\"uchteten 2016--2019, SOEP-Core v36 (Remote Edition), gewichtet.
    Lineare Wahrscheinlichkeitsmodelle; robuste, auf Personenebene geclusterte
    Standardfehler in Klammern. Bei standardisierten Variablen (std.) entspricht
    der Koeffizient dem Effekt einer Erh\"ohung um eine Standardabweichung.
    $^{+}$~p\,$<$\,0{,}1, $^{*}$~p\,$<$\,0{,}05, $^{**}$~p\,$<$\,0{,}01
    (zweiseitig). Alle 37 Koeffizienten- und Standardfehlerzellen stimmen im
    Signifikanzniveau mit \textcite[20]{Kanas:2023} \"uberein; die gr\"o\ss{}te
    Abweichung betr\"agt 0{,}28~p.p.
    \\[0.4ex]
    \textbf{Abweichung:} Die Kategorien des Asylstatus sind in
    der publizierten Tabelle falsch zugeordnet. Das Modell verwendet
    \texttt{ib1.asyl\_req\_stat}, und \texttt{00013\_data\_clean.do:1541}
    definiert Kategorie~1 als \emph{keine Entscheidung}; die Referenz ist somit
    \glqq keine Entscheidung\grqq{} und nicht \glqq anerkannt\grqq{}. Die
    Zuordnung oben ist die korrigierte und wurde gegen die Wertelabels der
    \texttt{mi estimate}-Ausgabe gepr\"uft. Inhaltlich bedeutet dies: sowohl
    anerkannte als auch abgelehnte Gefl\"uchtete arbeiten h\"aufiger als jene,
    deren Verfahren noch l\"auft.}
\end{table}

% Author  : Emre Suemer, 2026-08-13
%==============================================================================%

\makeatletter
\@ifundefined{NC@rewrite@Z}{%
  \newcolumntype{Z}{>{\raggedright\arraybackslash\hangindent=1.9em\hangafter=1}X}%
}{}
\makeatother

\begin{table}[htbp]
  \centering
  \caption{Wahrscheinlichkeit bezahlter Arbeit, marginale Effekte in Prozentpunkten}
  \label{tab:erwerbstaetigkeit}
  \footnotesize
  \setlength{\tabcolsep}{4pt}%
  \begin{tabularx}{\textwidth}{@{}Zrr@{}}
    \toprule
     & Modell 1.1 & Modell 1.2 \\
     & p.p. (SE) & p.p. (SE) \\
    \midrule
    \textbf{Kreisangebot an Integrationskursen, std.} & 2{,}98$^{*}$ (1{,}19) & 0{,}89 (1{,}78) \\
    Aufenthaltsdauer, in Monaten (/12) & 10{,}75$^{**}$ (1{,}60) & 10{,}99$^{**}$ (1{,}62) \\
    \quad $\times$ Kreisangebot an Integrationskursen &  & 0{,}81 (0{,}75) \\
    Bildungsjahre vor der Migration & 0{,}56$^{**}$ (0{,}19) & 0{,}56$^{**}$ (0{,}19) \\
    Erwerbserfahrung vor der Migration & 4{,}62$^{+}$ (2{,}39) & 4{,}57$^{+}$ (2{,}39) \\
    Weiblich & $-$15{,}83$^{**}$ (2{,}15) & $-$15{,}83$^{**}$ (2{,}15) \\
    Kinder unter 16 Jahren im Haushalt & $-$10{,}91$^{**}$ (2{,}15) & $-$10{,}90$^{**}$ (2{,}15) \\
    Einreise mit der Familie & 5{,}29$^{**}$ (1{,}96) & 5{,}33$^{**}$ (1{,}96) \\
    Unterst\"utzung durch Familie/Verwandte in DE & 4{,}10 (2{,}85) & 4{,}19 (2{,}85) \\
    Alter bei der Erstbefragung & $-$0{,}35$^{**}$ (0{,}09) & $-$0{,}35$^{**}$ (0{,}09) \\
    Traumatische Erfahrungen & $-$0{,}34 (2{,}31) & $-$0{,}36 (2{,}31) \\
    Status des Asylantrags (Ref.: keine Entscheidung) & & \\
    \quad -- Anerkannt & 5{,}20$^{*}$ (2{,}60) & 5{,}19$^{*}$ (2{,}60) \\
    \quad -- Abgelehnt & 8{,}24$^{**}$ (2{,}77) & 8{,}31$^{**}$ (2{,}76) \\
    Kreisanteil Ausl\"ander:innen, std. & $-$1{,}16 (2{,}79) & $-$1{,}12 (2{,}79) \\
    Kreisbev\"olkerungsdichte, std. & $-$0{,}89 (3{,}60) & $-$0{,}90 (3{,}59) \\
    Kreisarbeitslosenquote, std. & $-$4{,}45$^{+}$ (2{,}43) & $-$4{,}54$^{+}$ (2{,}44) \\
    Kreissorgen \"uber Zuwanderung, std. & 0{,}36 (1{,}20) & 0{,}34 (1{,}20) \\
    Kreisanteil CDU/CSU-W\"ahler:innen, std. & $-$1{,}92 (2{,}19) & $-$1{,}98 (2{,}20) \\
    Konstante & $-$2{,}63 (6{,}30) & $-$2{,}77 (6{,}29) \\
    \midrule
    \multicolumn{3}{@{}l}{\textit{Fixe Effekte}} \\
    \quad Erhebungsjahr & Ja & Ja \\
    \quad Erhebungsstichprobe & Ja & Ja \\
    \quad Bundesland & Ja & Ja \\
    \quad Herkunftsland (aggregiert) & Ja & Ja \\
    \midrule
    N Beobachtungen & 5.467 & 5.467 \\
    N Personen & 2.730 & 2.730 \\
    N Imputationen & 20 & 20 \\
    R\textsuperscript{2} (adj.) & 0{,}2135 & 0{,}2138 \\
    \bottomrule
  \end{tabularx}

  \vspace{0.75ex}
  \parbox{\textwidth}{\footnotesize
    \textit{Anmerkungen:} Eigene Berechnung, IAB-BAMF-SOEP-Befragung von
    Gefl\"uchteten 2016--2019, SOEP-Core v36 (Remote Edition), gewichtet.
    Lineare Wahrscheinlichkeitsmodelle; robuste, auf Personenebene geclusterte
    Standardfehler in Klammern. Bei standardisierten Variablen (std.) entspricht
    der Koeffizient dem Effekt einer Erh\"ohung um eine Standardabweichung.
    $^{+}$~p\,$<$\,0{,}1, $^{*}$~p\,$<$\,0{,}05, $^{**}$~p\,$<$\,0{,}01
    (zweiseitig). Alle 37 Koeffizienten- und Standardfehlerzellen stimmen im
    Signifikanzniveau mit \textcite[20]{Kanas:2023} \"uberein; die gr\"o\ss{}te
    Abweichung betr\"agt 0{,}28~p.p.
    \\[0.4ex]
    \textbf{Abweichung:} Die Kategorien des Asylstatus sind in
    der publizierten Tabelle falsch zugeordnet. Das Modell verwendet
    \texttt{ib1.asyl\_req\_stat}, und \texttt{00013\_data\_clean.do:1541}
    definiert Kategorie~1 als \emph{keine Entscheidung}; die Referenz ist somit
    \glqq keine Entscheidung\grqq{} und nicht \glqq anerkannt\grqq{}. Die
    Zuordnung oben ist die korrigierte und wurde gegen die Wertelabels der
    \texttt{mi estimate}-Ausgabe gepr\"uft. Inhaltlich bedeutet dies: sowohl
    anerkannte als auch abgelehnte Gefl\"uchtete arbeiten h\"aufiger als jene,
    deren Verfahren noch l\"auft.}
\end{table}

% Author  : Emre Suemer, 2026-08-13
%==============================================================================%

\makeatletter
\@ifundefined{NC@rewrite@Z}{%
  \newcolumntype{Z}{>{\raggedright\arraybackslash\hangindent=1.9em\hangafter=1}X}%
}{}
\makeatother

\begin{table}[htbp]
  \centering
  \caption{Wahrscheinlichkeit bezahlter Arbeit, marginale Effekte in Prozentpunkten}
  \label{tab:erwerbstaetigkeit}
  \footnotesize
  \setlength{\tabcolsep}{4pt}%
  \begin{tabularx}{\textwidth}{@{}Zrr@{}}
    \toprule
     & Modell 1.1 & Modell 1.2 \\
     & p.p. (SE) & p.p. (SE) \\
    \midrule
    \textbf{Kreisangebot an Integrationskursen, std.} & 2{,}98$^{*}$ (1{,}19) & 0{,}89 (1{,}78) \\
    Aufenthaltsdauer, in Monaten (/12) & 10{,}75$^{**}$ (1{,}60) & 10{,}99$^{**}$ (1{,}62) \\
    \quad $\times$ Kreisangebot an Integrationskursen &  & 0{,}81 (0{,}75) \\
    Bildungsjahre vor der Migration & 0{,}56$^{**}$ (0{,}19) & 0{,}56$^{**}$ (0{,}19) \\
    Erwerbserfahrung vor der Migration & 4{,}62$^{+}$ (2{,}39) & 4{,}57$^{+}$ (2{,}39) \\
    Weiblich & $-$15{,}83$^{**}$ (2{,}15) & $-$15{,}83$^{**}$ (2{,}15) \\
    Kinder unter 16 Jahren im Haushalt & $-$10{,}91$^{**}$ (2{,}15) & $-$10{,}90$^{**}$ (2{,}15) \\
    Einreise mit der Familie & 5{,}29$^{**}$ (1{,}96) & 5{,}33$^{**}$ (1{,}96) \\
    Unterst\"utzung durch Familie/Verwandte in DE & 4{,}10 (2{,}85) & 4{,}19 (2{,}85) \\
    Alter bei der Erstbefragung & $-$0{,}35$^{**}$ (0{,}09) & $-$0{,}35$^{**}$ (0{,}09) \\
    Traumatische Erfahrungen & $-$0{,}34 (2{,}31) & $-$0{,}36 (2{,}31) \\
    Status des Asylantrags (Ref.: keine Entscheidung) & & \\
    \quad -- Anerkannt & 5{,}20$^{*}$ (2{,}60) & 5{,}19$^{*}$ (2{,}60) \\
    \quad -- Abgelehnt & 8{,}24$^{**}$ (2{,}77) & 8{,}31$^{**}$ (2{,}76) \\
    Kreisanteil Ausl\"ander:innen, std. & $-$1{,}16 (2{,}79) & $-$1{,}12 (2{,}79) \\
    Kreisbev\"olkerungsdichte, std. & $-$0{,}89 (3{,}60) & $-$0{,}90 (3{,}59) \\
    Kreisarbeitslosenquote, std. & $-$4{,}45$^{+}$ (2{,}43) & $-$4{,}54$^{+}$ (2{,}44) \\
    Kreissorgen \"uber Zuwanderung, std. & 0{,}36 (1{,}20) & 0{,}34 (1{,}20) \\
    Kreisanteil CDU/CSU-W\"ahler:innen, std. & $-$1{,}92 (2{,}19) & $-$1{,}98 (2{,}20) \\
    Konstante & $-$2{,}63 (6{,}30) & $-$2{,}77 (6{,}29) \\
    \midrule
    \multicolumn{3}{@{}l}{\textit{Fixe Effekte}} \\
    \quad Erhebungsjahr & Ja & Ja \\
    \quad Erhebungsstichprobe & Ja & Ja \\
    \quad Bundesland & Ja & Ja \\
    \quad Herkunftsland (aggregiert) & Ja & Ja \\
    \midrule
    N Beobachtungen & 5.467 & 5.467 \\
    N Personen & 2.730 & 2.730 \\
    N Imputationen & 20 & 20 \\
    R\textsuperscript{2} (adj.) & 0{,}2135 & 0{,}2138 \\
    \bottomrule
  \end{tabularx}

  \vspace{0.75ex}
  \parbox{\textwidth}{\footnotesize
    \textit{Anmerkungen:} Eigene Berechnung, IAB-BAMF-SOEP-Befragung von
    Gefl\"uchteten 2016--2019, SOEP-Core v36 (Remote Edition), gewichtet.
    Lineare Wahrscheinlichkeitsmodelle; robuste, auf Personenebene geclusterte
    Standardfehler in Klammern. Bei standardisierten Variablen (std.) entspricht
    der Koeffizient dem Effekt einer Erh\"ohung um eine Standardabweichung.
    $^{+}$~p\,$<$\,0{,}1, $^{*}$~p\,$<$\,0{,}05, $^{**}$~p\,$<$\,0{,}01
    (zweiseitig). Alle 37 Koeffizienten- und Standardfehlerzellen stimmen im
    Signifikanzniveau mit \textcite[20]{Kanas:2023} \"uberein; die gr\"o\ss{}te
    Abweichung betr\"agt 0{,}28~p.p.
    \\[0.4ex]
    \textbf{Abweichung:} Die Kategorien des Asylstatus sind in
    der publizierten Tabelle falsch zugeordnet. Das Modell verwendet
    \texttt{ib1.asyl\_req\_stat}, und \texttt{00013\_data\_clean.do:1541}
    definiert Kategorie~1 als \emph{keine Entscheidung}; die Referenz ist somit
    \glqq keine Entscheidung\grqq{} und nicht \glqq anerkannt\grqq{}. Die
    Zuordnung oben ist die korrigierte und wurde gegen die Wertelabels der
    \texttt{mi estimate}-Ausgabe gepr\"uft. Inhaltlich bedeutet dies: sowohl
    anerkannte als auch abgelehnte Gefl\"uchtete arbeiten h\"aufiger als jene,
    deren Verfahren noch l\"auft.}
\end{table}
